\section{A worked gate analysis with adaptive selection}\label{sec:C25:workedgate}
Consider a system that proposes a code change, obtains a test verdict, asks a model critic for approval, and publishes the accepted artifact for review. There are at least four distinct latent events: the patch satisfies the requirement, the test suite covers the relevant behavior, the critic interprets the evidence correctly, and the published artifact is the one that was checked. A failure at any stage can produce a misleading acceptance. Calling the arrangement a verification pipeline does not remove those events from the analysis.

Begin with the artifact-binding obligation. The test runner records digest $d_1$, while the publication tool is about to send digest $d_2$. If $d_1\ne d_2$, the gate should reject or rerun regardless of the critic's confidence. This is a deterministic check over recorded identities. If the digests match, the system has established that evidence and artifact correspond, but it has not established that the tests cover the specification. The distinction allows an engineering invariant to be enforced without claiming an unavailable semantic oracle.

Next consider a candidate that passes all visible tests by special-casing the tested examples. A model critic may notice the shortcut, but it may also infer that the behavior is acceptable because the benchmark passes. A requirement-grounded witness can target an untested input implied by the prompt. If the witness fails on the candidate and passes on the reference, the remaining question is whether its expectation actually follows from the task. An independent reviewer should be shown that question without being encouraged to treat the reference's incidental choices as requirements.

Now add adaptive search. The agent sees a rejection and revises the patch. The resulting candidates are not independent draws from a fixed distribution: each can be optimized against prior feedback. A false-acceptance rate measured before this interaction may no longer apply. To obtain a union bound over $n$ attempts, one sufficient condition is a bound $f_t$ on the conditional probability of false acceptance at attempt $t$ given the entire previous search history. Then the probability of any false acceptance is at most $\min(1,\sum_t f_t)$. Marginal calibration on unrelated samples does not establish these conditional bounds.

The implementation consequence is to retain the search history and gate version. If a candidate generator is allowed to query the gate repeatedly, the evaluation must expose that access and include it in the budget. A final accepted artifact should identify how many candidates were considered and whether the gate's feedback was visible during revision. Otherwise two systems with the same final test suite can have different effective opportunities to exploit it.

A holdout gate can reduce direct adaptation, but it introduces its own governance problem. If researchers repeatedly inspect holdout failures to improve the system, the holdout gradually becomes development data. Rotating tasks or using a genuinely untouched confirmation set can help, provided the task distribution remains relevant. A smaller untouched sample may yield wider uncertainty than a large heavily reused benchmark. Scientific reporting should preserve that uncertainty rather than prefer the larger but compromised number.

Threshold selection also depends on consequence. A gate used to prioritize human review has a different loss function from a gate that authorizes external publication. Let $c_{FA}>0$ be the cost of false acceptance and $c_{FR}>0$ the cost of false rejection. For $q=\Pr(Y=1\mid O)$ conditional on the information $O$ actually available to that decision, accepting has expected error cost $c_{FA}(1-q)$ and rejecting has expected error cost $c_{FR}q$ under a simplified two-action model. Acceptance is favored when $q\geq c_{FA}/(c_{FA}+c_{FR})$, with either decision permissible at equality. An estimated score requires validation in the corresponding decision population; marginal calibration on unrelated candidates does not supply this posterior when additional covariates change the costs or error distribution. The formula does not make either probability or cost available automatically.

An unresolved or review action changes that decision problem. Review can acquire additional evidence at cost, and the reviewer is also fallible. A rational policy compares expected residual loss after review with immediate acceptance or rejection. In practice these quantities are difficult to estimate, so applications often use hard rules for consequential effects and probabilistic gates for lower-impact filtering. The rules should be justified by the action boundary, not by an arbitrary global confidence threshold.

This analysis transfers to internal model monitors. Suppose a sparse feature is associated with a concerning behavior and triggers a warning. The feature's activation can be a useful signal without being a complete causal explanation. Calibration depends on model revision, context distribution, and intervention history. If the deployment policy changes behavior in response to the monitor, the observed data distribution changes too. A monitor evaluated only before intervention can therefore become miscalibrated after it is embedded in a control loop.

The discriminating experiment has separate discovery, calibration, adaptive-attack, and confirmation phases. Discovery identifies candidate signals or tests. Calibration estimates their behavior on a declared population. Adaptive evaluation allows a bounded optimizer to search for accepted failures. Confirmation evaluates the final fixed configuration on untouched cases. Report how candidates, tasks, and model families overlap between phases. This makes the independence claim inspectable rather than relying on the label held out.

Finally, preserve the distinction between a falsifying witness and a comprehensive safety claim. One replayable counterexample can refute a universal assertion about a gate. The absence of a counterexample in a finite search cannot establish that no accepted failure exists. It can support a bounded empirical statement about the search procedure, budget, and task population. That statement is valuable when accurately scoped. Inflating it into certification makes the evidence less useful precisely because it hides the operating conditions under which the gate was actually examined.

The ROME implementation makes this separation concrete. Its published causal-tracing experiment exposes a patching routine and hidden-flow calculation, while the model-editing module exposes the function that applies ROME updates\cite{meng2022rome}. These are distinct execution paths with distinct claims: restoring selected activations tests mediation under an intervention; applying an update tests whether a weight change produces the desired association and preserves specified neighboring behavior. A reproduction ledger should not mark both claims reproduced because only the editing demo ran. The sparse-autoencoder implementation likewise separates encoding and decoding in its Autoencoder class\cite{gao2024sae}. Reconstructing activations through those functions establishes numerical behavior of the learned dictionary, while interpreting a feature and evaluating its downstream causal effect requires additional experiments. Record the model checkpoint, selected layer, activation normalization, feature selection rule, intervention magnitude and metric. These fields identify where two apparently contradictory analyses may actually differ. Code availability therefore strengthens inspectability without collapsing the distinction between a reachable function, an executed experiment and a supported scientific conclusion.

In particular, the inspected ROME source contains \path{trace_with_patch}, \path{calculate_hidden_flow}, \path{apply_rome_to_model} and \path{execute_rome}. Their presence was checked; no execution result is asserted here. Pinning these callables to a revision and running the associated controls is the implementation step required before an independent reproduction claim.
